\section{6교시 (75분) — 실습 ② 종료 보고서·미해결 이관·교훈·편익 원장}

% =====================================================================
% 6교시 — 실습 ②
% =====================================================================
\begin{frame}{6교시 — 이 교시에서 하는 것}
\begin{itemize}
\item 5교시의 세 문서(⑱·⑲·⑳)를 \textbf{손으로} 쓴다 — 워크북 §3\til{}5
\item 세 과제, 한 실 — 이관 6번(나영선) \ra{} 교훈 적용 대상 \ra{} 원장 BR-05
\item ⑱ 종료 보고서 최종 4축 + 이관 목록 5건 이상 (25분)
\item ⑲ 교훈 5건 — 여섯 칸, 적용 대상·증빙까지 (20분)
\item ⑳ 편익 원장 B-01\til{}B-06 × 6칸, 측정 책임자 실명 (25분)
\end{itemize}
\keymsg{종료는 선언이 아니라 증명이다 — 세 문서의 모든 숫자·이름은 앞 문서의 어느 행에서 왔는지 추적된다}
\note{Day3 6교시와 같은 운영 — 개인 작성 후 조별 대조, 강사는 순회. 다른 점은 세 과제가 "프로젝트가 끝난 뒤에도 남는 것"을 적는다는 것이다 — 이관 목록은 P1이 하지 않은 것, 교훈은 다음 프로젝트가 바꿀 것, 원장은 운영이 잴 것. 세 문서는 한 실로 이어진다: ⑱ 이관 6번(입고 검사기록 디지털화, 나영선, 2027-09 이사회)이 ⑲ 개정 요청 ①(BC 편익 표에 분모 칸)의 근거가 되고 ⑳ B-02 분모 B의 목표 설정 조건(BR-05)으로 다시 나타난다. 수강생에게 워크북 §3.7·§4.7·§5.7 기본 과제와 §3.5·§4.5·§5.5 완성 예시본을 펴게 한다. 시간이 부족하면 ⑲는 3건, ⑳은 B-02·B-04·B-07 세 항목으로 줄이되 6칸은 다 채운다. 예상 소요 3분.}
\end{frame}

\begin{frame}{과제 ⑱ — 워크북 §3: 종료 보고서 최종 4축 + 이관 목록}
\begin{itemize}
\item 시점 2027-05-28 G4 — 워크북 §3.7 기본 과제 상황(SAC 항목 7 미충족)
\item 최종 4축 — 범위 FP 12,585 / 일정 G4 = 기준선 완료일 / 원가 AC 151.95 / 품질 SAC 11+1
\item 원가는 3층을 다 적는다 — 168 승인 / 156 한도 / 146.82 PMB \ra{} 여유 4.05 · VAC $-$5.13
\item 이관 목록 5건 이상 — 출처·소유자 실명·기한·승인/승인 전·수신 조직
\item 헌장 §14 ① 두 경로 중 택일 — "미해소 0" / "해소 기한·책임자 서명"
\end{itemize}
\keymsg{원가 한 층만 적으면 "절감"이 되고, 세 층을 적으면 "16.05 미사용·여유 4.05·초과 5.13"이 동시에 참이 된다}
\note{워크북 §3.3 항목 1·2·6과 §3.5 예시본 항목 2·6을 쓴다. 기본 과제는 SAC 항목 7(SLA/OLA/UC)이 UC ② 재협상 결렬로 미충족인 상황 — 06-18(허용범위 끝)까지 G4를 미루는 대안과 05-28 조건부 종료의 비용(PMO·수행사 잔류·하자보수 개시 지연·유보금)을 비교하고 항목 1·6·7·8을 다시 쓴다. 최종 4축은 그림 4-1(3층 + EAC 이력 9호 150.86 \ra{} 11호 152.40 \ra{} 17호 151.95)의 값을 그대로 인용한다 — 종료 시점에 새로 만든 숫자는 없어야 한다. 이관 목록은 예시본 10건(승인 6·승인 전 4) 중 5건 이상을 자기 문장으로 다시 쓰되, 미충족 항목 7이 "승인"인지 "승인 전"인지를 반드시 판정하게 한다(UC 예산 증액 결정자가 누구인가). 예상 소요 25분.}
\end{frame}

\begin{frame}{미해결 사항 이관 10건 — 출처에서 수신 조직까지}
\fig[0.6]{../그림/PM_Day4/fig_4_4_followon_map.pdf}
\figcap{그림 4-4. 출처(SAC 조건부·CR 조건·SLA 강등·편익 전제) \ra{} 항목·소유자 실명·기한 \ra{} 수신 조직. 승인 6 실선·승인 전 4 파선. 6번(나영선·2027-09 이사회)이 ⑳ BR-05로 이어진다}
\keymsg{이관 목록에 "이 프로젝트의 편익이 실현되려면 이 프로젝트가 하지 않은 것"이 한 줄이라도 있는가}
\note{워크북 §3.5 항목 6 표의 10건을 지도로 옮긴 것. 승인 전 4건 — 3번(가맹 마감 22:00 전환, 한동석, 2027-09 CAB), 6번(입고 검사기록 디지털화, 나영선, 2027-09 이사회), 7번(SAP ECC 유지보수 종료 대응, 박세연, 별도 과제), 10번(PMO 잔류 예산, 프로그램 PMO장, 06-01 운영위) — 은 무엇이 승인되어야 하는지(누가·어느 회의체)를 적었기 때문에 종료와 함께 사라지지 않는다. 특히 6번은 P1 범위 밖(종이·엑셀 검사기록)이지만 재고 정확도 98.5\% 목표의 전제이므로 적었다 — 2025-11 이사회가 34억 SI 견적을 반려한 이력까지. 수강생의 목록에서 소유자가 부서명이거나 기한이 "추후"인 행을 찾아 표시하게 한다. 예상 소요 5분.}
\end{frame}

\begin{frame}{과제 ⑲ — 워크북 §4: 교훈 5건, 여섯 칸}
\begin{itemize}
\item 상황 · 원인 · 교훈 · \textbf{적용 대상} · 소유자 · \textbf{증빙} — 여섯 칸을 다 채운다
\item 상황은 숫자로 — 12영업일 · 0.42억 · +38분 · D$-$2 발견
\item 원인은 사람이 아니라 양식·규칙·구조 — "그 칸이 없었다"
\item 5건 중 잘된 것 1건 이상(L-08 동시 기록 · L-15 분리 컷오버 연결)
\item 5건을 개정 요청 3건(규정·내부통제 지침·계약 문안)으로 묶는다
\end{itemize}
\keymsg{"적용 대상"이 없는 교훈은 감상이다 — 문서명과 칸까지 적히고 그 문서의 소유자가 접수해야 교훈이다}
\note{워크북 §4.3 항목 1·2와 §4.5 예시본(15건)을 쓴다. 기본 과제는 헌장 §14 중단 조건이 2026-11-27 G2에서 발동했다고 가정한 premature closure의 교훈 10건인데, 시간상 이 교시에서는 정상 종료 교훈 5건으로 줄여 진행하고 중단 시나리오는 과제로 남긴다. 교재 §5.3의 다섯 건(L-02 대체 결정권자, L-06 감리 귀속 기준, L-08 동시 기록, L-09 컷오버 창 앞 사전 결재, L-15 분리 컷오버의 연결)을 본으로 삼되, 수강생은 자기 문장으로 다시 쓴다. 채점의 핵심은 원인 칸에 사람 이름이 있는지, 소유자가 PM인지(PM은 곧 없어진다), 증빙이 워크북 산출물의 항목·행까지 적혔는지 세 가지다. 잘된 것의 교훈이 없으면 다음 프로젝트는 이번 성공을 다시 발명해야 한다. 예상 소요 20분.}
\end{frame}

\begin{frame}{교훈이 사라지는 경로 — Syllk 6층과 여섯 칸}
\fig[0.6]{../그림/PM_Day4/fig_5_1_syllk.pdf}
\figcap{그림 5-1. 사람 3층 / 시스템 3층, 6개월 소멸선. ④ 적용 대상·⑤ 소유자 칸이 process·infrastructure로 가는 문이다 — 개정 요청 3건(07·08·09월)이 그 문을 통과한다}
\keymsg{교훈은 아카이브에서 죽는다 — 개정 요청서는 등록부의 부속이 아니라 결과물이다}
\note{교재 §5.1 Williams(2008)와 Duffield \& Whitty의 Syllk. 사람 층(learning·culture·social)에 머문 교훈은 팀 해산과 함께 6개월 안에 소멸한다는 것이 그림의 소멸선이다. 시스템 층(technology·process·infrastructure)으로 넘어가는 유일한 문이 여섯 칸의 ④ 적용 대상과 ⑤ 소유자다. 온담 P1은 06-04에 개정 요청 3건을 발송했다 — ① 프로젝트 관리 규정·양식집(프로그램 PMO장, 2027-07 운영위) ② 재무본부 내부통제 지침(재경팀장 정미란, 컷오버 창 앞 사전 결재 목록) ③ 구매 계약 표준 문안(강현도·최영주, 감리 귀속·동시 지연·대기 기록). 수강생의 5건이 이 셋 중 어디로 가는지 표시하게 한다. 예상 소요 4분.}
\end{frame}

\begin{frame}{과제 ⑳ — 워크북 §5: 편익 원장 B-01\til{}B-06 × 6칸}
\begin{itemize}
\item 6칸 — 측정식 · 기준값 · 목표값 · 측정 시점 · \textbf{측정 책임자 실명} · 미달 시 조치
\item 측정식은 코드 수준 — 분자·분모·산출 시스템(RPT-FRN-03)·주기
\item 기준값에는 반드시 분모와 측정 시점 — 91.3\% (순환실사 SKU, 05-21)
\item 미달 시 조치는 원인 분기 — output / outcome / utilisation
\item 직접 편익 118억, 온담 귀속 99.6억 — 두 기준 실현률을 따로 적는다
\end{itemize}
\keymsg{프로젝트가 끝난 뒤 이 원장을 들고 있는 조직과 사람의 이름이 적혀 있는가 — 그것이 편익 원장의 첫 조건이다}
\note{워크북 §5.3 항목 1·2·5와 §5.5 예시본을 쓴다. 온담 P1의 6항목 — B-04 폐기율 43억·B-05 배달수수료 37억·B-07 인건비율 38억·B-02 재고 정확도·B-03 결품률·B-08 리드타임(참고 2 — B-01 온라인 비중은 B-05로 환산, B-06 처리량은 P2 이관). 과제명은 B-01\til{}B-06 여섯 칸 표이나 항목 번호는 예시본의 B-ID를 그대로 쓰게 한다. 측정 책임자는 편익 소유자 실명 + 측정 실무 직책 — 오정근 5·한동석 1·재무본부 판정, ITO 서명표(항목 5)의 세 열과 같은 사람이어야 한다. 미달 시 조치는 "원인 분석"이 아니라 output이면 하자보수·유지보수, outcome이면 업무 변화 소유자, utilisation이면 P3 채택 조치로 분기한다. 시간이 없으면 B-02·B-04·B-07 세 항목만 6칸으로. 예상 소요 25분.}
\end{frame}

\begin{frame}{ITO 사슬과 책무 분리 — ⑤ 칸은 서명란이다}
\fig[0.6]{../그림/PM_Day4/fig_5_2_ito_ledger.pdf}
\figcap{그림 5-2. Input \ra{} Transform \ra{} Output ǀ Utilisation \ra{} Outcome. 서명자 3열 — output P1 PM·이재현 / utilisation 매장운영팀장·강현도·나영선·박준영 / outcome 오정근 5·한동석 1·재무본부 판정 — 과 원장 6칸(B-04 행)}
\keymsg{output 소유자(PM)는 종료 후 편익 판정에서 빠진다 — 편익을 잰 사람이 편익을 만든 사람이면 판정이 아니다}
\note{교재 §5.4 Zwikael \& Smyrk의 ITO 모델. 책임 분리선의 왼쪽(Input·Transform·Output)은 프로젝트, 오른쪽(Utilisation·Outcome)은 운영이다. 원장 ⑤ 칸(측정 책임자)은 outcome 열의 사람이어야 하며 PM이 들어가면 안 된다 — PM은 G4 다음 날 없고, 있더라도 자기 output을 스스로 판정하게 된다. 수강생의 원장에서 측정 책임자가 "영업본부"이거나 PM이면 오류로 표시한다. utilisation 열이 비어 있는 항목(채택이 안 되면 편익이 안 나는 항목 — B-05 앱 전환)은 P3 채택 조치와 연결해야 한다. 예상 소요 4분.}
\end{frame}

\begin{frame}{작성 요령 ①②③ — 이관 한 줄 · 적용 대상 · 기준값 확정}
\begin{tbl}[\scriptsize]{L{2.2cm}L{4.6cm}L{4.6cm}}
요령 & 규칙 & 온담 P1의 한 줄 \\ \midrule
① 이관 한 줄의 5요소 & 출처 · 소유자 실명 · 기한 · 승인/승인 전(누가·어느 회의체) · 수신 조직 & 6번: 편익 원장 B-02 전제 · 나영선 · 2027-09 이사회 · 승인 전 · 생산본부 \\
② 교훈의 적용 대상 & 세 층 중 하나 이상 — 다음 프로젝트 양식 / 사내 표준 조항 / 계약 문안 · 문서명·칸까지 & L-09: 내부통제 지침 "컷오버 창 앞 4주 사전 결재 목록" 조항 신설 · 재경팀장 \\
③ 기준값 확정 & 승인본이 전제한 분모로 판정 · 다른 분모는 병기 · \textbf{이유를 적는다} & A 91.3\%(순환실사 SKU, 승인본 전제) 확정 · B 84.6\% 병기, 목표 M+12 설정 \\
\end{tbl}
\keymsg{91.3인가 84.6인가 — 어느 쪽이든 값 옆에 "왜 이 분모인가"를 적지 않으면 G5에서 판정권 다툼이 된다}
\note{워크북 §3.3 항목 6, §4.3 항목 1, §5.3 항목 2의 "어떻게 적는가"를 한 표로 압축했다. ①은 다섯 요소 중 하나라도 비면 그 행은 종료와 함께 소멸한다 — 특히 "승인 전"은 무엇이 승인되어야 하는지를 적어야 다음 회의체 안건이 된다. ②는 "전사"·"향후 프로젝트"가 적용 대상이 아니라는 뜻이다 — 개정될 문서의 이름과 조항·칸, 그 문서의 소유 부서를 적는다. ③은 Day3 §5.6 정정(91.3)이 편익 판정권의 문제였음을 다시 확인하는 자리다 — 승인본이 약속한 분모로 판정하는 것이 정직하고, 확장 분모의 목표는 P1이 하지 않은 전제(검사기록 디지털화, 이관 6번) 위에 세울 수 없다. 유리한 분모로 바꾸거나 둘 중 하나를 지우는 것이 흔한 오류. 예상 소요 6분.}
\end{frame}

\begin{frame}{분모 정치 — A 91.3 \ra{} 98.5 와 B 84.6, 판정과 목표는 다른 문서다}
\fig[0.6]{../그림/PM_Day4/fig_5_3_denominator.pdf}
\figcap{그림 5-3. A 91.3\ra{}98.5(7.2\%p) vs B 84.6(13.9\%p 필요, 1.93배). 비순환 SKU 31\%(이천·안성 원료·반제품, 검사기록 34억 반려). 판정 A 서명 05-21 \ra{} B 목표 조건(이관 6번) \ra{} 가결/부결(BR-05) 분기}
\keymsg{두 분모를 다 적고 판정 분모를 명시한다 — 하나를 지우는 순간 나머지 하나가 정치가 된다}
\note{교재 §5.6과 워크북 ⑳ 항목 2. 분모 B로 판정하면 목표까지 13.9\%p가 필요해 A의 1.93배 — 같은 프로젝트가 성공에서 실패로 바뀐다. 그래서 온담은 승인본 전제인 A로 판정 기준값을 확정하고(05-21 서명), B는 병기하되 목표는 M+12 리뷰(2028-05)에서 검사기록 과제의 이사회 결정(2027-09, 이관 6번)에 따라 설정한다 — 가결이면 B 목표 설정, 부결이면 B는 참고 지표로 남긴다(원장 BR-05). 수강생 과제 ⑳의 B-02 행에 이 분기가 적혔는지 확인한다. §5.7 기본 과제(G5 실적 B 88.1\%)는 이 분기를 실제로 실행하는 연습이다. 예상 소요 4분.}
\end{frame}

\begin{frame}{흔한 오류 5 · 예시본과 템플릿 · 시간 배분}
\begin{columns}[T]
\begin{column}{0.52\textwidth}
\subhead{흔한 오류 5}
\begin{itemize}
\item 원가 실적을 한 층(168 승인)과만 비교 — "절감 16.05"
\item 이관 소유자가 부서명, 기한이 "추후", 승인 전을 승인처럼
\item 교훈의 원인이 사람, 적용 대상이 "전사", 소유자가 PM
\item 기준값에 분모·측정 시점 없음, 유리한 분모로 교체
\item 측정 책임자가 "영업본부" · 미달 시 조치가 "원인 분석"
\end{itemize}
\end{column}
\begin{column}{0.46\textwidth}
\subhead{예시본·템플릿}
\begin{itemize}
\item ⑱ §3.4 템플릿 · §3.5 예시본 · §3.6 해설
\item ⑲ §4.4 템플릿 · §4.5 예시본 15건
\item ⑳ §5.4 템플릿 · §5.5 예시본 · 그림 5-2·5-3
\end{itemize}
\subhead{시간 배분 · 심화}
\begin{itemize}
\item ⑱ 25 / ⑲ 20 / ⑳ 25 / 대조 5
\item 심화(3\til{}7년차) — 자기 회사 종료 보고서·교훈·편익 표를 §3.7·§4.7·§5.7 심화 과제로 센다
\end{itemize}
\end{column}
\end{columns}
\keymsg{셋 중 둘 이상이 없으면 그 보고서는 종료를 선언한 문서이지 종료를 증명한 문서가 아니다}
\note{순회하며 다섯 오류를 표시하게 한다 — 대부분의 오류는 칸이 비어 있는 것이 아니라 칸이 이름 대신 부서명·"추후"·"전사"·"원인 분석"으로 채워진 것이다. 예시본은 베끼는 것이 아니라 대조하는 것 — 수강생의 값이 예시본과 다르면 워크북이 정본이지만, 다른 이유를 적을 수 있으면 그것이 더 좋은 답이다. 심화 과제 셋은 회사에 돌아가 하는 것: ① 원가가 어느 층과 비교됐는가·다른 층으로 바꾸면 절감이 초과로 바뀌는가 ② 이관 목록에 실명·승인 전 표시가 몇 개인가 ③ 지체상금 정산에 근거 문장이 있는가 / 교훈 3건의 여섯 칸 중 적용 대상·소유자를 채울 수 있는가·실제 개정된 양식이 하나라도 있는가 / 승인 문서 편익 표에서 종료 후 실제 측정된 항목이 몇 개인가(0이면 2020 감사 지적이 그 회사에도 해당). 예상 소요 5분.}
\end{frame}

\begin{frame}{6교시 핵심 정리 — 프로젝트가 끝난 뒤에 남는 세 문서}
\begin{itemize}
\item ⑱ 종료 보고서 — 모든 숫자는 성과보고·변경 로그·등록부의 행에서 온다
\item 이관 목록 — 5요소가 다 있어야 종료와 함께 소멸하지 않는다
\item ⑲ 교훈 — 상황은 숫자, 원인은 구조, 적용 대상은 문서명과 칸
\item ⑳ 원장 — 6칸, 측정 책임자는 outcome 열의 실명, 판정 분모를 명시
\item 세 문서는 한 실 — 이관 6번 \ra{} 개정 요청 ① \ra{} BR-05, 끊기면 편익은 승인 시점에만 존재한 숫자
\end{itemize}
\keymsg{끝났다고 선언한 사람·받았다고 서명한 사람·승인한 사람이 세 사람인가 — 그리고 3년 뒤 이 원장을 들고 있을 사람의 이름이 있는가}
\note{정리 3분. 세 점검 질문으로 닫는다 — ⑱: 종료 보고서의 모든 숫자가 앞 문서의 어느 행에서 왔는지 추적되는가(추적되지 않는 숫자는 종료 시점에 만들어진 숫자다) / ⑲: 원인 칸에 사람 이름이 없고 적용 대상의 문서 소유자가 개정 요청을 접수했는가 / ⑳: 편익 항목마다 output·outcome·utilisation 소유자가 세 사람이고 output 소유자가 판정에서 빠져 있는가. 마무리 30분(제6장, 그림 6-1 20종 사슬)으로 넘어가며 ⑳에서 ①로 되돌아가는 여섯 가닥(측정 책임자·분모·손익분기·AC·지체상금 0·PoNR)을 예고한다. 예상 소요 3분.}
\end{frame}
